\section{Conclusion}
\label{sec:conclusion}

This work combines teacher-assistant distillation, low-bit quantization-aware
training, and FINN dataflow compilation to obtain a deployable Ultra96
accelerator for the four-class ultra-widefield fundus task. The intermediate
ResNet-18 assistant was needed to avoid the capacity-gap loss observed when the
large teacher was distilled directly into the tiny student. Six-bit QAT preserved
the student's accuracy while reducing the arithmetic precision enough for FPGA
deployment.

The hardware exploration was ultimately limited by on-chip memory rather than
compute. Moving selected FIFOs from block RAM to LUTRAM made the full 6-bit
\texttt{test\_resnet} fit, but it left almost no block-RAM headroom for extra
folding. The reduced-width variant narrowed the late channel widths, freed memory,
and enabled one additional folding point. The resulting operating points span
clock scaling and folding: at the same 187.5\,MHz clock, throughput rose from
6.36 images/s for the full model to 12.71 images/s for the reduced-width model,
with board accuracy decreasing only from 81.20 to 80.45. The final
reduced-width accelerator is bit-exact with its software reference and reaches
79.89 weighted F1 and 75.62 macro F1.

The remaining limitations are the ones exposed by the design space itself.
Four-bit quantization collapses accuracy for this small network, 8-bit hardware
does not fit the Ultra96 resource budget, and the final deployment remains
memory-bound. The main clinical weakness is still recall on the minority lesion
classes, especially melanoma and CHRPE. Future work should therefore improve
FIFO sizing and folding selection in the FINN flow, explore architectures with
better minority-class recall at the same memory cost, and evaluate larger FPGA
devices able to host higher-resolution variants.
